% !TeX root = ../main.tex

% ==================================================
% 1.5. Python File I/O
% ==================================================


\subsection{Python File I/O}

가장 밑바닥인 \zsh{open()}과 file object가 정확히 무엇인가부터 시작한다.

\subsubsection{\zsh{open()}은 무엇을 하는가?}

가장 기본적인 코드는:

\begin{dhkcode}
f = open("data.txt", "r")
\end{dhkcode}

이다.

흔히 이것을 그냥 "\zsh{data.txt}를 연다"라고 배우는데,

객체 관점에서 보면 조금 더 정확하게 이해할 수 있다.

\begin{dhkbranch}[in=8]
node main :: "data.txt"
arrow main :: 경로를 open()에 전달
node main :: open("data.txt", "r")
arrow main :: OS에 파일 접근 요청
node main :: file object 생성
arrow main
node main :: f → file object
\end{dhkbranch}

즉,

\begin{dhkcode}
f = open("data.txt", "r")
\end{dhkcode}

에서 \zsh{f}가 파일 자체인 것은 아니다.

\zsh{f}는 \textbf{열려 있는 파일과 상호작용하기 위한 Python의 file object}를 가리키는 이름이다.

이건 지금까지 배운 객체 모델과 완전히 동일하다.

\begin{dhkdir}
f
> file object
>> 실제 파일과 연결
\end{dhkdir}

그래서 다음과 같은 메서드를 사용할 수 있다.

\begin{dhkcode}
f.read()
f.readline()
f.write(...)
f.close()
\end{dhkcode}

\subsubsection{\zsh{"r"}은 무엇인가?}

두 번째 인자는 \textbf{mode}이다.

\begin{dhkcode}
open("data.txt", "r")
\end{dhkcode}

여기서 \zsh{"r"}은 \textbf{read}.

즉:

\begin{dhksmry}
이 파일을 읽기 위한 방식으로 열어라.
\end{dhksmry}

라는 뜻이다.

대표적인 mode부터 보면:

\begin{dhksmry}
    "r" read   읽기
    "w" write  쓰기
    "a" append 뒤에 추가
\end{dhksmry}

예를 들어 파일 내용이:

\begin{dhkcode}
epoch 1
epoch 2
\end{dhkcode}

일 때,

\begin{dhkcode}
f = open("log.txt", "a")
f.write("\nepoch 3\n")
\end{dhkcode}

하면 기존 내용을 유지하면서 뒤에 추가해서:

\begin{dhkzsh}
epoch 1
epoch 2
epoch 3
\end{dhkzsh}

가 된다.

반면:

\begin{dhkcode}
f = open("log.txt", "w")
\end{dhkcode}

의 \zsh{"w"}는 \textbf{기존 파일이 있으면 내용을 비우고 새로 쓰기 시작할 수 있기 때문에 주의해야 한다}.

\subsubsection{파일을 읽어보자}

\zsh{data.txt}:

\begin{dhkcode}
apple
banana
orange
\end{dhkcode}

Python:

\begin{dhkcode}
f = open("data.txt", "r")
text = f.read()
print(text)
f.close()
\end{dhkcode}

\zsh{f.read()}는 파일에서 내용을 읽어 Python 객체로 가져온다.

여기서는 text mode이므로 대략:

\begin{dhkbranch}[in=8]
node main :: disk
node main :: data.txt
arrow main :: read
node main :: Python process
node main :: text → "apple\nbanana\norange"
\end{dhkbranch}

처럼 \zsh{str}객체를 얻게 된다.

그래서:

\begin{dhkcode}
type(text)
\end{dhkcode}

는:

\begin{dhkzsh}
str
\end{dhkzsh}

이다.

\subsubsection{그런데 왜 \zsh{close()}가 필요한가?}

이게 File I/O에서 중요한 부분이다.

\begin{dhkcode}
f = open(...)
\end{dhkcode}

하면 Python은 단순히 파일 이름을 기억하는 게 아니라 운영체제의 파일 자원을 사용하게 된다.

그래서 사용이 끝났으면:

\begin{dhkcode}
f.close()
\end{dhkcode}

로 닫아줘야 한다.

문제는 이런 코드다:

\begin{dhkcode}
f = open("data.txt", "r")

text = f.read()

# 여기에서 에러 발생
some_function()

f.close()
\end{dhkcode}

\zsh{some_function()}에서 예외가 발생하면?

\begin{dhksmry}[centered]
open

\da

read

\da

에러 발생

\da

프로그램 흐름 중단

\da

f.close()까지 못 감
\end{dhksmry}

이런 자원 관리 문제 때문에 Python에서 일반적으로 

직접 \zsh{open()} $\rightarrow$ \zsh{close()}를 관리하기보다는 \zsh{with}를 사용한다.

\subsubsection{\zsh{with open(...) as f}}

실무에서 흔히 보는 형태가 이것이다.

\begin{dhkcode}
with open("data.txt", "r") as f:
    text = f.read()

print(text)
\end{dhkcode}

핵심은:

\begin{dhksmry}[centered]
with 블록 진입

\da

파일 열기

\da

f \ra file object

\da

f.read()

\da

with 블록 종료

\da

파일 자동 정리/close
\end{dhksmry}

중간에 예외가 발생하더라도 \zsh{with}가 자원 정리를 처리할 수 있다.

그래서 보통:

\begin{dhkcode}
f = open(...)

...

f.close()
\end{dhkcode}

보다:

\begin{dhkcode}
with open(...) as f:
    ...
\end{dhkcode}

를 선호한다.

\subsubsection{그런데 \zsh{with}는 파일 전용 문법이 아니다}

이 부분이 중요하다.

처음 배우면:

\begin{dhksmry}[centered]
with \ra 파일 열 때 사용하는 문법
\end{dhksmry}

처럼 외우기 쉬운데 그렇지 않다.

\zsh{with}는 \textbf{context manager}라는 더 일반적인 Python 개념과 관련된다.

큰 그림은:

\begin{dhksmry}[centered]
어떤 자원을 사용하기 전에 준비

\da

사용

\da

사용이 끝나면 반드시 정리
\end{dhksmry}

같은 패턴을 안전하게 다루는 문법이다.

파일이 대표적인 예일 뿐이다.

\begin{dhkcode}
with open(...) as f:
    ...
\end{dhkcode}

를 개념적으로 읽으면:

\begin{dhknote}
\begin{dhksmry}
이 context에 진입하면서 자원을 준비하고, 
블록 안에서 사용한 뒤 context를 빠져나갈 때 정리한다.
\end{dhksmry}
\end{dhknote}

라고 볼 수 있다.

나중에 class를 더 배우면 \zsh{__enter__, __exit__}과 연결된다.

\subsubsection{\zsh{Path}와도 바로 연결된다}

앞에서 배운 \zsh{pathlib}을 그대로 사용할 수 있다.

\begin{dhkcode}
from pathlib import Path

path = Path("data/train.txt")

with path.open("r") as f:
    text = f.read()
\end{dhkcode}

즉:

\begin{dhkcode}
open("data/train.txt", "r")
\end{dhkcode}

과

\begin{dhkcode}
path = Path("data/train.txt")
path.open("r")
\end{dhkcode}

둘 다 파일을 여는 방식이다.

그리고 단순히 텍스트 전체를 읽는 경우에는:

\begin{dhkcode}
text = path.read_text()
\end{dhkcode}

처럼 더 간단한 API도 있다.

\subsubsection{여기서 첫 번째로 잡아야 할 핵심}

지금은 메서드를 외우는 것보다 이 구조가 중요하다.

\begin{dhkbranch}[in=8]
node main :: file system
arrow main :: path
node main :: "data/train.txt" 또는 Path(...)
arrow main
node main :: open()
arrow main
node main :: file object
branchline main -> read :: read()
stop read
branchline main -> rl :: readline()
stop rl
branchline main -> write :: write()
stop write
handoff main -> close :: close()
stop close
\end{dhkbranch}

그리고:

\begin{dhksmry}[centered]
name f
\da
file object
\da
OS의 열린 파일 지원
\da
실제 filesystem의 파일
\end{dhksmry}

이렇게 \textbf{Python 객체 세계와 실제 filesystem 사이를 연결하는 객체가 file object}이다.

다음은 \zsh{read(), readline(), readlines()}의 차이와

왜 \textbf{file object가 iterator이기도 하는지}로 넘어간다.

\subsubsection{\zsh{read()} - 파일 전체를 한 번에 읽기}

\zsh{data.txt}가 다음과 같다고 하자.

\begin{dhkcode}
apple
banana
orange
\end{dhkcode}

\begin{dhkcode}
with open("data.txt", "r") as f:
    text = f.read()

print(text)
\end{dhkcode}

\zsh{text}에는 하나의 \zsh{str}객체가 들어온다.

\begin{dhkzsh}
'apple\nbanana\norange'
\end{dhkzsh}

구조적으로는:

\begin{dhkbranch}
node main :: file
arrow main :: f.read()
node main :: 하나의 str 객체
node main :: "apple\nbanana\norange"
\end{dhkbranch}

즉, \zsh{read()}는 \textbf{전체 파일 $\rightarrow$ 하나의 문자열}이라고 생각하면 된다.

그래서 큰 파일에서는 문제가 생길 수 있다.

\begin{dhkbranch}
node main ::10GB file
arrow main :: read()
node main :: 10GB에 가까운 데이터를 메모리에 한 번에 올리려고 함
\end{dhkbranch}

\subsubsection{\zsh{readline()} - 한 줄씩 읽기}

이번에는:

\begin{dhkcode}
with open("data.txt", "r") as f:
    line1 = f.readline()
    line2 = f.readline()

print(line1)
print(line2)
\end{dhkcode}

대략:

\begin{dhkzsh}
line1 == "apple\n"
line2 == "banana\n"
\end{dhkzsh}

이 된다.

여기서 중요한 현상이 하나 있다.

첫 번째:

\begin{dhkcode}
f.readline()
\end{dhkcode}

은 \zsh{apple}을 읽었는데,

두 번째:

\begin{dhkcode}
f.readline()
\end{dhkcode}

은 또 \zsh{apple}을 읽는 게 아니라 \zsh{banana}를 읽었다.

왜 그럴까?

\subsubsection{file object에는 현재 읽고 있는 위치가 있다}

file object는 내부적으로 \textbf{현재 파일의 어느 위치까지 읽었는지}를 관리한다.

개념적으로:

\begin{dhkdir}
apple\nbanana\norange
>현재위치
\end{dhkdir}

첫 번째 \zsh{readline()}:

\begin{dhkdir}
apple\n@banana\norange
>현재 위치
\end{dhkdir}

두 번째 \zsh{readline()}:

\begin{dhkdir}
apple\nbanana\n@orange
>현재위치
\end{dhkdir}

처럼 진행된다.

이 위치를 보통 \textbf{file position/file pointer}라고 부른다고 생각하자.

실제로 확인할 수도 있다.

\begin{dhkcode}
with open("data.txt", "r") as f:
    print(f.tell())

    f.readline()

    print(f.tell())
\end{dhkcode}

\zsh{tell()}은 현재 위치를 알려준다.

그리고:

\begin{dhkcode}
f.seek(0)
\end{dhkcode}

처럼 \zsh{seek()}을 사용하면 위치로 이동할 수도 있다.

\subsubsection{그래서 \zsh{read()}도 두 번 하면?}

이제 결과를 예측할 수 있다.

\begin{dhkcode}
with open("data.txt", "r") as f:
    a = f.read()
    b = f.read()

print(a)
print(b)
\end{dhkcode}

a 에는:

\begin{dhkzsh}
apple
banana
orange
\end{dhkzsh}

가 들어가지만 b 는:

\begin{dhkcode}
''
\end{dhkcode}

빈 문자열이 된다.

왜냐하면 첫 번째 \zsh{read()}가:

\begin{dhksmry}
    [apple banana orange]
    \ua

    f.read()

    [apple banana orange]
                        \ua
                        EOF
\end{dhksmry}

파일 끝까지 위치를 이동시켰기 때문이다.

\zsh{EOF}는 \textbf{End Of File}이다.

따라서 두 번째 \zsh{read()}에는 더 이상 읽을 것이 없다.

\subsubsection{\zsh{readlines()} - 모든 줄을 리스트로}

이번에는:

\begin{dhkcode}
with open("data.txt", "r") as f:
    lines = f.readlines()

print(lines)
\end{dhkcode}

결과는 대략:

\begin{dhkzsh}
['apple\n', 'banana\n', 'orange']
\end{dhkzsh}

이다.

즉:

\begin{dhksmry}[centered]
read()
\da
\zsh{"apple\nbanana\norange"}
하나의 str


readline()
\da
\zsh{"apple\n"}
한 줄 str


readlines()
\da
\zsh{['apple\n', 'banana\n', 'orange']}
list[str]
\end{dhksmry}

여기까지만 보면 \zsh{readlines()}가 편해보인다.

하지만 역시 파일 전체의 줄들을 리스트로 만들어야 하기에 큰 파일에서는 메모리를 사용하게 된다.

\subsubsection{그런데 Python에서는 파일을 그냥 \zsh{for}에 넣을 수 있다}

여기서 \zsh{iterator}가 등장한다.

\begin{dhkcode}
with open("data.txt", "r") as f:
    for line in f:
        print(line)
\end{dhkcode}

이 가능하다.

왜?

\textbf{file object는 iterato이기 때문이다}.

그래서 개념적으로:

\begin{dhkcode}
iterator = iter(f)

line1= next(iterator)
line2= next(iterator)
line3= next(iterator)
\end{dhkcode}

가 가능하다.

실제로:

\begin{dhkcode}
with open("data.txt", "r") as f:
    print(iter(f) is f)
\end{dhkcode}

는 일반적인 Python text file object에서:

\begin{dhkzsh}
True
\end{dhkzsh}

가 된다.

즉 file object 자신이 \zsh{iterator}이다.

\subsubsection{iterator protocol이 그대로 나온다}

이터레이터 구조가:

\begin{dhksmry}[centered]
iterable

\da

iter()

\da

iterator

\da

next()

\da

다음 값

\da

...

\da

StopIteration
\end{dhksmry}

이다.

file object에서는:

\begin{dhksmry}[centered]
file object

\da

iterator

\da

next(f)

\da

한 줄

\da

next(f)

\da

다음 한 줄

\da

...

\da

EOF

\da

StopIteration
\end{dhksmry}

이다.

그래서:

\begin{dhkcode}
with open("data.txt", "r") as f:
    print(next(f))
    print(next(f))
\end{dhkcode}

하면:

\begin{dhkzsh}
apple
banana
\end{dhkzsh}

순서로 읽힌다.

\subsubsection{\zsh{for line in f}가 큰 파일에서 중요한 이유}

예를 들어 학습 데이터가:

\begin{dhkzsh}
10GB trianing_data.txt
\end{dhkzsh}

라고 해보자.

이걸:

\begin{dhkcode}
with open("training_data.txt") as f:
    text = f.read()
\end{dhkcode}

하면 전체 파일을 한꺼번에 읽으려고 한다.

\begin{dhksmry}[centered]
10GB file
\da
read()
\da
거대한 str
\da
memory
\end{dhksmry}

반면:

\begin{dhkcode}
with open("training_data.txt") as f:
    for line in f:
        process(line)
\end{dhkcode}

이면 개념적으로:

\begin{dhkdir}
disk
> line 1 -> memory -> 처리
> line 2 -> memory -> 처리
> line 3 -> memory -> 처리
> ...
\end{dhkdir}

처럼 \textbf{전체 파일을 하나의 거대한 Python 문자열/리스트로 만들 필요가 없다}.

물론 실제 OS/Python I/O는 성능을 위해 buffering을 사용하기 때문에

디스크에서 정확히 물리적인 한 줄씩 읽는다는 뜻은 아니다.

하지만 \textbf{Python 코드 관점에서는 한 줄씩 순차적으로 소비할 수 있다}는 것이 핵심이다.

\subsubsection{여기서 아주 중요한 연결 하나}

위에서:

\begin{dhkcode}
with open("data.txt") as f:
    for line in f:
        ...
\end{dhkcode}

에서 \zsh{f}가 \zsh{iterator}라고 했다.

그러면:

\begin{dhkcode}
with open("data.txt") as f:
    for line in f:
        print(line)
    
    for line in f:
        print(line)
\end{dhkcode}

두 번재 \zsh{for}문에서는 어떻게 될까?

\textbf{아무것도 출력되지 않는다}.

첫 번째 \zsh{for}가 이미:

\begin{dhksmry}[centered]
f
\da
apple
\da
banana
\da
orange
\da
EOF
\end{dhksmry}

까지 소비했기 때문이다.

이것이:

\begin{dhknote}
\begin{dhksmry}[centered]
소진된 iterator는 다시 처음부터 시작하지 않는다.
\end{dhksmry}
\end{dhknote}

와 정확히 같은 현상이다.

처음부터 다시 읽으려면 예를 들어:

\begin{dhkcode}
f.seek(0)
\end{dhkcode}

으로 파일 위치를 처음으로 되돌려야 한다.

\subsubsection{지금 단계에서 핵심 정리}

\begin{dhktable}
|함수|내용|
|:---|:---|
|read()|남은 내용을 하나의 str로 읽음|
|readline()|현재 위치부터 한 줄을 str로 읽음|
|readlines()|남은 줄들을 list[str]로 읽음|
|for line in f|file iterator를 순차 소비|
|next(f)|다음 한 줄|
| tell()|현재 file position 확인|
|seek(...)|file position 이동|
\end{dhktable}

그리고 가장 중요한 개념은:

\begin{dhknote}
\begin{dhksmry}[centered]
파일 읽기는 단순히 파일에서 값을 복사해오는 것이 아니라,

file object가 현재 읽기 위치라는 상태를 가지고 

그 상태가 진행되는 과정이다.
\end{dhksmry}
\end{dhknote}

이것을 이해하면 \zsh{read()} 두 번 했는데 왜 두 번째가 비어있는지,

\zsh{readline()}이 왜 다음 줄로 넘어가는지,

파일을 왜 \zsh{for}문으로 순회할 수 있는지가 전부 하나의 원리로 연결된다

다음 단계에서는 쓰기(w, a, write)와 파일이 실제로 언제 만들어지고 언제 내용이 사라지는지 살펴본다.

\subsubsection{\zsh{"w"} - 파일을 쓰기 모드로 열기}

\begin{dhkcode}
with open("result.txt", "w") as f:
    f.write("epoch 1")
\end{dhkcode}

실행하면 \zsh{result.txt}에:

\begin{dhkzsh}
epoch 1
\end{dhkzsh}

이 저장된다.

그런데 \zsh{"w"}에는 중요한 성질이 두 가지 있다.

\textbf{파일이 없응면 새로 만들고, 파일이 이미 있으면 기존 내용을 비운다}.

예를 들어 원래:

\begin{dhkcode}
epoch 1
epoch 2
epoch 3
\end{dhkcode}

였는데:

\begin{dhkcode}
with open("result.txt", "w") as f:
    f.write("new training")
\end{dhkcode}

하면 결과는:

\begin{dhkzsh}
new training
\end{dhkzsh}

이다.

즉 \zsh{"w"}는 기존 내용 뒤에 이어 쓰는 모드가 아니다.

그래서 \zsh{"w"}는 잘못 사용하면 기존 파일 내용을 날릴 수 있다.

\subsubsection{\zsh{"a"} - append}

기존 내용을 보존하면서 \textbf{끝에 추가}하려면:

\begin{dhkcode}
with open("result.txt", "a") as f:
    f.write("epoch 4\n")
\end{dhkcode}

을 사용한다.

원래:

\begin{dhkzsh}
epoch 1
epoch 2
epoch 3
\end{dhkzsh}

이었다면:

\begin{dhkzsh}
epoch 1
epoch 2
epoch 3
epoch 4
\end{dhkzsh}

가 된다.

ML 학습 로그를 계속 누적한다고 생각하면 직관적이다.

\begin{dhkcode}
with open("training.log", "a") as f:
    f.write("epoch=10 loss=0.231\n")
\end{dhkcode}

\subsubsection{\zsh{wirte()}는 자동으로 줄바꿈하지 않는다}

이것도 자주 실수하는 부분이다.

\begin{dhkcode}
with open("result.txt", "w") as f:
    f.write("apple")
    f.write("banana")
    f.write("orange")
\end{dhkcode}

결과:

\begin{dhkzsh}
applebananaorange
\end{dhkzsh}

\zsh{print()}처럼 자동 줄바꿈을 해주는 것이 아니다.

따라서:

\begin{dhkcode}
f.write("apple\n")
f.write("banana\n")
f.write("orange\n")
\end{dhkcode}

처럼 줄바꿈 이스케이프 문자를 넣어야 한다.

\subsubsection{\zsh{write()}는 무엇을 반환할까?}

재미있는 부분인데:

\begin{dhkcode}
n = f.write("hello")
\end{dhkcode}

\zsh{n}에는 보통:

\begin{dhkzsh}
5
\end{dhkzsh}

가 들어간다.

즉 \zsh{write()}는 \textbf{작성한 문자 수}를 반환한다.

그래서 \zsh{write()}의 반환값은 \zsh{None}이 아니라는 것도 알아두면 좋다.

\subsubsection{여러 문자열을 쓸 때 \zsh{writelines()}}

예를 들어:

\begin{dhkcode}
lines = [
    "epoch 1\n",
    "epoch 2\n",
    "epoch 3\n"
]

with open("result.txt", "w") as f:
    f.writelines(lines)
\end{dhkcode}

결과:

\begin{dhkzsh}
epoch 1
epoch 2
epoch 3
\end{dhkzsh}

여기서도 \zsh{writelines()}가 자동으로 \zsh{\n}을 붙여주는 것은 아니다.

즉:

\begin{dhkcode}
["apple", "banana", "orange"]
\end{dhkcode}

를 \zsh{writelines()} 하면:

\begin{dhkzsh}
applebananaorange
\end{dhkzsh}

가 된다.

\subsubsection{그런데 \zsh{write()}하면 즉시 디스크에 기록되는가?}

여기서 buffering이 등장한다.

처음에 보기엔:

\begin{dhksmry}[centered]
f.write(...)
\da
disk
\end{dhksmry}

라고 생각하기 쉽지만 실제로는 보통 더 복잡하다.

개념적으로는:

\begin{dhksmry}[centered]
Python program

f.write("hello")
\da
buffer
\da
적절한 시점
\da
OS
\da
filesystem
\end{dhksmry}

중간에 buffer를 두는 것이다.

왜 굳이 그럴까?

디스크나 OS에 아주 작은 데이터를 쓸 때마다 계속 요청하는 것보다

어느 정도 모아서 처리하는 것이 효율적이기 때문이다.

예를 들어:

\begin{dhkcode}
f.write("a")
f.write("b")
f.write("c")
f.write("d")
\end{dhkcode}

라고 해서 반드시 물리 저장장치에:

\begin{dhkzsh}
a 기록
b 기록
c 기록
d 기록
\end{dhkzsh}

와 같은 식으로 각각 독립적으로 기록되는 것은 아니다.

\subsubsection{\zsh{flush()}}

그래서 등장하는 것이:

\begin{dhkcode}
f.flush()
\end{dhkcode}

이다.

개념적으로는:

\begin{dhksmry}[centered]
Python 쪽에 버퍼링된 출력 데이터를 아래 계층으로 밀어내라.
\end{dhksmry}

정도로 이해하면 된다.

예를 들어 장시간 학습 프로그램이:

\begin{dhkcode}
f = open("training.log", "w")
for epoch in range(1000):
    ...
    f.write(f"epoch={epoch}, loss={loss}\n")
\end{dhkcode}

처럼 돌아간다고 해보자.

프로그램이 몇 시간 동안 실행될 수도 있다.

이때 로그를 외부에서 빨리 확인해야 한다면

buffering 때문에 원하는 시점에 내용이 안 보이는 상황이 있을 수 있다.

그럴 때:

\begin{dhkcode}
f.write(...)
f.flush()
\end{dhkcode}

처럼 명시적으로 flush할 수도 있다.

\subsubsection{그러면 \zsh{close()}와 \zsh{flush()}는 무슨 관계인가?}

일반적으로 파일을 정상적으로 \zsh{close()}하면 필요한 \zsh{flush()}도 수행된다.

즉:

\begin{dhkcode}
f.write(...)
f.close()
\end{dhkcode}

에서는 close 과정에서 버퍼에 남아 있는 데이터를 내보내고 파일 자원을 정리한다.

그리고 앞에서 배운:

\begin{dhkcode}
with open("result.txt", "w") as f:
    f.write("hello")
\end{dhkcode}

도 블록을 정상적으로 빠져나오면서 file object가 정리되니까

일반적인 파일 작업에서는 직접 \zsh{flush()}를 계속 호출할 필요가 없다.

그래서 \zsh{with}가 다시 중요해진다.



\begin{dhkbranch}[in=8]
node main :: with open(...) as f:
branch main -> write1 :: write
stop write1
branch main -> write2 :: write
stop write2
branch main -> write3 :: write
stop write3
node main :: with 종료
arrow main
node main :: 필요한 flush
arrow main
node main :: close
arrow main
node main :: 자원 정리
\end{dhkbranch}

\subsubsection{단, \zsh{flush()}가 물리 디스크에 영구 저장 완료는 아니다}

이건 지금 단계에서는 구분만 해둔다.

\begin{dhkcode}
f.flush()
\end{dhkcode}

했다고 해서 반드시 그 순간 \textbf{SSD의 물리 셀까지 영구적으로 기록되었다}는 의미는 아니다.

대략 여러 계층이 존재할 수 있다.

\begin{dhksmry}[centered]
Python

\da

Python buffering

\da

OS

\da

OS filesystem/cache

\da

storage device
\end{dhksmry}

\zsh{flush()}는 기본적으로 Python의 버퍼를 아래 계층으로 넘기는 것과 관련된 것이지,

모든 하드웨어 계층까지 물리 기록 완료를 보장한다는 의미는 아니다.

정말 그런 수준의 보장이 필요한 경우라면

\zsh{os.fsync()}같은 별도의 개념이 등장하는데, 지금은 살펴볼 필요가 없다.

\subsubsection{\zsh{"r", "w", "a"}를 하나로 정리하면}

\begin{dhksmry}
    "r"
    읽기
    파일이 없으면 \ra 오류
    기존 내용 보존


    "w"
    쓰기
    파일이 없으면 \ra 생성
    파일이 있으면 \ra 기존 내용 제거
    
    
    "a"
    추가 쓰기
    파일이 없으면 \ra 생성
    파일이 있으면 \ra 기존 내용 보존 후 끝에 추가
\end{dhksmry}

이후에:

\begin{dhksmry}
    "x" 새 파일 생성, 이미 존재하면 오류


    "+" 읽기 + 쓰기


    "b" binary mode


    "t" text mode
\end{dhksmry}

까지 붙으면 파일 mode 전체가 완성된다.

\begin{dhksmry}[centered]
    번외로, bytes와 encoding, decoding에 대한 내용

    "r"
    \da
    읽겠다

    text mode
    \da
    bytes를 문자로 해석하겠다



    encoding="utf-8"
    \da
    그 bytes를 어떤 규칙으로 문자로 바꿀 것인가


    파일의 bytes
    \da
    UTF-8 규칙으로 decode
    \da
    Python str
    "안녕하세요"


    text mode
    \ra Python이 encoding/decoding을 대신 처리
    \ra Python 코드에서는 str 사용

    binary mode
    \ra encoding/decoding하지 않음
    \ra Python 코드에서는 bytes 사용


    encoding은 정확히는 str \ra bytes 로 바꾸는 규칙
    decoding은 그 반대인 bytes \ra str 로 해석하는 규칙


    encoding='utf-8' 은 파일 타입을 지정하는 게 아니라
    문자(str)와 파일에 저장되는 바이트(bytes) 사이를 
    어떤 규칙으로 변환할지 지정한다
    cp949 도 똑같은 역할인데 변환 규칙이 UTF-8과 다른 것


    bytes 는 저수준 프로그래밍 언어가 아니라 
    컴퓨터가 파일·네트워크 등에서 다루는 바이트 데이터 표현
    머릿속에서는 간단하게:
    사람이 다루는 문자(str) \ra encode \ra bytes \ra decode \ra 문자(str)
\end{dhksmry}

저수준의 File I/O는 이 정도로 끝내고,

실제로 Python/ML/DL에서 자주만나는 구조화된 파일 형식으로 넘어간다.

\subsubsection{JSON이 왜 필요한가?}

지금까지는 이런 텍스트를 저장했다.

\begin{dhkzsh}
epoch=10
batch_size = 32
learning_rate = 0.001
\end{dhkzsh}

사람은 이해할 수 있지만 Python 입장에서는 그냥 하나의 문자열이다.

ML 설정을 생각해보자.

\begin{dhkzsh}
config = {
    "epochs": 10,
    "batch_size": 32,
    "learning_rate": 0.001,
    "hidden_sizes": [512, 256, 128]
}
\end{dhkzsh}

이런 \textbf{구조화된 데이터}를 파일로 저장하고 다시 불러오고 싶을 수 있다.

JSON은 이런 데이터를 표현하기 위한 대표적인 형식이다.

\zsh{config.json}:

\begin{dhkzsh}
{
    "epochs": 10,
    "batch_size": 32,
    "learning_rate": 0.001,
    "hidden_sizes": [512, 256, 128]
}
\end{dhkzsh}

Python Dictionary와 상당히 비슷하게 생겼다.

하지만 중요한 차이가 존재한다.

\begin{dhknote}
\begin{dhksmry}[centered]
JSON 자체는 Python dictionary가 아니다.
\end{dhksmry}
\end{dhknote}

그냥 JSON이라는 데이터 형식이다.

\subsubsection{JSON 파일 $\rightarrow$ Python 객체}

Python에서는 표준 라이브러리 \zsh{json}을 사용한다.

\begin{dhkcode}
import json

with open("config.json", "r", encoding="utf-8") as f:
    config = json.load(f)
\end{dhkcode}

그러면:

\begin{dhkcode}
print(type(config))
\end{dhkcode}

결과:

\begin{dhkzsh}
<class 'dict'>
\end{dhkzsh}

그리고:

\begin{dhkcode}
print(config['learning_rate'])
\end{dhkcode}

\begin{dhkzsh}
0.01
\end{dhkzsh}

처럼 사용할 수 있다.

흐름을 보면:

\begin{dhksmry}[centered]
config.json

\da

file object f

\da

json.load(f)

\da

Python 객체

\da

config

\da

dict
\end{dhksmry}

즉 \zsh{json.load()}는 JSON 형식의 데이터를 읽어서 적절한 Python 객체로 변환해준다.

\subsubsection{JSON과 Python 타입은 서로 대응한다}

대표적으로:

\begin{dhktable}
|JSON|Python|
|:---|:---|
|object|dict|
|array|list|
|string|str|
|number|int/float|
|true|True|
|false|False|
|null|None|
\end{dhktable}

그래서 JSON:

\begin{dhkzsh}
{
    "use_gpu": true,
    "dropout": null,
    "layers": [512, 256]
}
\end{dhkzsh}

을 읽으면 Python에서는 대략:

\begin{dhkzsh}
{
    "use_gpu": True,
    "dropout": None,
    "layers": [512, 256]
}
\end{dhkzsh}

가 된다.

\subsubsection{반대로 Python 객체 $\rightarrow$ JSON 파일}

이번에는:

\begin{dhkzsh}
config = {
    "epochs": 20,
    "batch_size": 64,
    "learning_rate": 0.001
}
\end{dhkzsh}

를 저장해보자.

\begin{dhkcode}
import json

with open("config.json", "w", encoding="utf-8") as f:
    json.dump(config, f)
\end{dhkcode}

방향이 반대다.

\begin{dhksmry}[centered]
Python dict

\da

json.dump()

\da

JSON 형식

\da

config.json
\end{dhksmry}

그래서 이름도 꽤 직관적이다.

\begin{dhksmry}
    load
    파일에서 Python으로 불러옴

    dump
    Python 객체를 파일로 내보냄
\end{dhksmry}

\subsubsection{\zsh{load()}와 \zsh{loads()}는 다르다}

여기서 Python에서 상당히 헛갈리는 이름이 등장한다.

\begin{dhkcode}
json.load(...)
json.loads(...)
\end{dhkcode}

s 하나 차이인데 의미가 다르다.

\textbf{json.load()}

파일 객체에서 읽는다.

\begin{dhkcode}
with open("config.json") as f:
    config = json.load(f)
\end{dhkcode}

\begin{dhksmry}[centered]
file object
\da
json.load()
\da
Python object
\end{dhksmry}


\textbf{json.loads()}

이미 가지고 있는 \textbf{JSON 문자열}을 Python 객체로 변환한다.

\begin{dhkzsh}
text = '{"epochs": 10, "batch_size": 32}'
config = json.loads(text)
\end{dhkzsh}

여기서는 파일이 없다.

\begin{dhksmry}[centered]
JSON str
\da
json.loads()
\da
Python dict
\end{dhksmry}

여기서 \zsh{loads}의 s는 \textbf{string}이라고 기억하면 편하다.

\subsubsection{반대쪽에도 \zsh{dump()}와 \zsh{dumps()}가 있다}

\begin{dhkcode}
json.dump(...)
json.dumps(...)
\end{dhkcode}

\textbf{dump()}:

\begin{dhksmry}[centered]
Python object
\da
JSON
\da
file object
\end{dhksmry}

\textbf{dumps()}:

\begin{dhksmry}[centered]
Python object
\da
JSON str
\end{dhksmry}

예를 들어:

\begin{dhkcode}
config = {
    "epochs": 10,
    "batch_size": 32
}

text = json.dumps(config)
print(type(text))
\end{dhkcode}

결과:

\begin{dhkzsh}
<class 'str'>
\end{dhkzsh}

즉 전체를 하나로 정리하면:

\begin{dhktable}
||파일|문자열|
|:---|:---|:---|
|읽기|json.load()|json.loads()|
||↓|↓|
||Python object|Python object|
|||
|쓰기|json.dump()|json.dumps()|
||↑|↑|
||Python object|Python object|
\end{dhktable}

\subsubsection{ML/DL에서는 어디에 쓰일까?}

예를 들어 학습 설정을:

\begin{dhkzsh}
{
    "model_name": "distilbert-base-uncased",
    "epochs": 3,
    "batch_size": 16,
    "learning_rate": 0.00002
}
\end{dhkzsh}

로 저장하고:

\begin{dhkcode}
import json

with open("config.json", encoding="utf-8") as f:
    config = json.load(f)

model_name = config["model_name"]
epochs = config["epochs"]
\end{dhkcode}

처럼 사용할 수 있다.

또 학습 결과를:

\begin{dhkcode}
metrics = {
    "accuracy": 0.91,
    "f1": 0.89,
    "loss": 0.21
}

with open("metric.json", "w", encoding="utf-8") as f:
    json.dump(metrics, f, indent=4)
\end{dhkcode}

처럼 저장할 수도 있다.

\zsh{indent=4}를 사용하면 사람이 읽기 좋게 들여쓰기해서 저장해준다.

\subsubsection{그런데 여기서 \zsh{argparse}와 연결된다}

우리가 이전에 배운:

\begin{dhkzsh}
python train.py \
    --epochs 10 \
    --batch-size 32 \
    --learning-rate 0.001
\end{dhkzsh}

도 결국 \textbf{프로그램 외부에서 설정을 넣는 방법}이었다.

JSON도 비슷한 목적에 사용할 수 있다.

\begin{dhkzsh}
python train.py --config config.json
\end{dhkzsh}

그리고:

\begin{dhkcode}
import argparse
import json
from pathlib import Path

parser = argparse.Argumentparser()
parser.add_argument(
    "--config",
    type=Path,
    required=True
)

args = parser.parse_args()

with args.config.open(encoding='utf-8') as f:
    config = json.load(f)

print(config['learning_rate'])
\end{dhkcode}

이제 지금까지 따로 배웠던 것들이:

\begin{dhksmry}[centered]
Terminal
\da
argparse
\da
Path
\da
file object
\da
JSON
\da
dict
\da
ML training configuration
\end{dhksmry}

으로 하나로 연결된다.

이게 앞에서:

\begin{dhksmry}[centered]
argparse \ra pathlib \ra filesystem \ra File I/O    
\end{dhksmry}

를 순서대로 공부한 이유이기도 하다.

다음은 CSV를 Python 표준 \zsh{csv} 모듈 관점에서 짧게 보고,

왜 ML/Data Science에서는 보통 그것보다

\zsh{pandas.read_csv()}를 쓰는지 연결해서 볼 것이다.

\subsubsection{CSV는 사실 텍스트 파일이다}

예를 들어 \zsh{scores.csv}:

\begin{dhkzsh}
name,math,english
kim,90,85
lee,80,92
park,95,88
\end{dhkzsh}

CSV는 \textbf{Comma-Separated Values}이다.

특별한 Python 객체가 저장되어 있는 게 아니라

기본적으로 이런 \textbf{텍스트 데이터}이다.

그래서 원한다면:

\begin{dhkcode}
with open("scores.csv", "r", encoding="utf-8") as f:
    print(f.read())
\end{dhkcode}

로도 읽을 수 있다.

다만 우리가 원하는 건 보통 문자열 전체가 아니라:

\begin{dhktable}
|name|math|english|
|:---|:---|:---|
|kim|90|85|
|lee|80|92|
|park|95|88|
\end{dhktable}

같은 \textbf{표 구조}이다.

그래서 CSV 형식을 해석하는 도구가 필요하다.

\subsubsection{Python에는 표준 \zsh{csv}모듈이 있다}

\begin{dhkcode}
import csv

with open("scores.csv", "r", encoding="utf-8") as f:
    reader = csv.reader(f)

    for row in reader:
        print(row)
\end{dhkcode}

결과는:

\begin{dhkzsh}
['name', 'math', 'english']
['kim', '90', '85']
['lee', '80', '92']
['park', '95', '88']
\end{dhkzsh}

여기서 중요한 부분이 하나 보인다

\begin{dhkzsh}
'90'
\end{dhkzsh}

이지:

\begin{dhkzsh}
90
\end{dhkzsh}

이 아니다.

CSV 자체는 텍스트이기 때문에 \zsh{csv.reader()}가 기본적으로 읽어온 필드들도 문자열로 다뤄진다.

필요하면 우리가:

\begin{dhkcode}
math = int(row[1])
\end{dhkcode}

처럼 변환할 수 있다.

\subsubsection{\zsh{csv.reader(f)}와 \zsh{iterator}}

여기서 배운 내용이 연결된다.

\begin{dhkcode}
reader = csv.reader(f)
\end{dhkcode}

\zsh{reader}도 순차적으로 row를 제공하는 iterable/iterator 계열 객체다.

그래서:

\begin{dhkcode}
for row in reader:
    ...
\end{dhkcode}

가 자연스럽게 가능하다.

흐름은:

\begin{dhksmry}[centered]
file object
\da
csv.reader
\da
row
\da
row
\da
row
\da
...
\end{dhksmry}

즉 Python의 \zsh{iterator}모델이 여기서도 계속 나타난다.

\subsubsection{그런데 ML/Data Science에서는 \zsh{Pandas}를 훨씬 많이 쓴다}

\begin{dhkcode}
import pandas as pd

df = pd.read_csv("scores.csv")
\end{dhkcode}

가 등장한다.

이 경우 Pandas가 상당히 많은 일을 대신 처리한다.

\begin{dhksmry}[centered]
CSV file
\da
파일 읽기
\da
CSV parsing
\da
header 처리
\da
dtype 추론
\da
결측값 처리
\da
DataFrame 생성
\end{dhksmry}

그래서:

\begin{dhkcode}
df.info()
df.describe()
df["math"].mean()
\end{dhkcode}

같은 데이터 분석 작업으로 바로 넘어갈 수 있다.

\subsubsection{CSV에서 실제로 알아둘 옵션}

ML/Data Science 기준으로는 이 정도면 충분하다.

\begin{dhkcode}
df = pd.read_csv(
    "data.csv",
    encoding='utf-8'
)
\end{dhkcode}

한글 파일이면 상황에 따라:

\begin{dhkcode}
df = pd.read_csv(
    "data.csv",
    encoding='cp949'
)
\end{dhkcode}

를 만날 수 있다.

쉼표가 아닌 다른 구분자를 사용하는 파일도 있다.

예를 들어:

\begin{dhkzsh}
name|math|english
kim|90|85
lee|80|92
\end{dhkzsh}

이면:

\begin{dhkcode}
df = pd.read_csv(
    "data.csv",
    sep="|"
)
\end{dhkcode}

처럼 delimiter를 지정할 수 있다.

\zsh{sep} = separator라고 보면 된다.

\subsubsection{DataFrame $\rightarrow$ CSV도 가능하다}

반대 방향:

\begin{dhkcode}
df.to_csv("result.csv")
\end{dhkcode}

그런데 Pandas에서 흔히:

\begin{dhkcode}
df.to_csv(
    "result.csv",
    index=False,
    encoding="utf-8"
)
\end{dhkcode}

를 많이 본다.

\zsh{index=False}는 DataFrame의 index를 별도 CSV열로 저장하지 않겠다는 뜻이다.

예를 들어:

\begin{dhksmry}
        name score
    0   kim    90
    1   lee    80
\end{dhksmry}

를 그냥 저장하면 index까지 포함될 수 있는데,

\begin{dhkzsh}
index = False
\end{dhkzsh}

로 하면:

\begin{dhkzsh}
name,score
kim,90
lee,80
\end{dhkzsh}

처럼 실제 데이터 열만 저장할 수 있다.

\subsubsection{그러면 CSV는 이 정도면 끝이다}

기억할 것은:

\begin{dhksmry}[centered]
CSV
\da
텍스트 기반의 표 형식


Python 기본 기능
\da
csv.reader / csv.writer


Data Science / ML
\da
보통 Pandas 이용


pd.read_csv()
\da
DataFrame


df.to_csv()
\da
CSV
\end{dhksmry}

이 정도면 충분하다.

\subsubsection{이제 Binary I/O}

\textbf{Text mode와 Binary mode}

지금까지:

\begin{dhkcode}
open("data.txt", "r")
\end{dhkcode}

라고 썼다.

사실 기본적으로 text mode라는 의미가 포함되어 있다.

명시하면:

\begin{dhkcode}
open("data.txt", "rt")
\end{dhkcode}

라고도 할 수 있다.

\zsh{t} = text

반면:

\begin{dhkcode}
open("image.jpg", "rb")
\end{dhkcode}

에서는:

\begin{dhksmry}[centered]
    r = read
    b = binary
\end{dhksmry}

즉 \textbf{binary read}이다.

\subsubsection{가장 중요한 차이는 반환되는 Python 객체}

\textbf{Text}:

\begin{dhkcode}
with open("data.txt", "r", encoding="utf-8") as f:
    data = f.read()

print(type(data))
\end{dhkcode}

\begin{dhkzsh}
str
\end{dhkzsh}

\textbf{Binary}:

\begin{dhkcode}
with open("image.jpg", "rb") as f:
    data = f.read()

print(type(data))
\end{dhkcode}

\begin{dhkzsh}
bytes
\end{dhkzsh}

즉:

\begin{dhksmry}[centered]
text mode

file bytes
\da
decoding
\da
Python str



binary mode

file bytes
\da
그대로
\da
Python bytes
\end{dhksmry}

이게 핵심이다.

\subsubsection{Binary에서는 encoding을 지정하지 않는다}

예를 들어:

\begin{dhkcode}
open("image.jpg", "rb", encoding="utf-8")
\end{dhkcode}

처럼 생각하면 \textbf{안된다}.

Binary mode에서는:

\begin{dhksmry}[centered]
bytes
\da
그대로 사용
\end{dhksmry}

하기 때문이다.

encoding은:

\begin{dhksmry}
bytes \lra str
\end{dhksmry}

변환 규칙이다.

Binary mode는 그 변환 자체를 하지 않는다.

\subsubsection{그런데 ML/DL에서는 보통 직접 안 한다}

예를 들어 NumPy array를 저장한다고 하자.

우리가 직접:

\begin{dhkcode}
with open("array.npy", "wb") as f:
    ...
\end{dhkcode}

하면서 NumPy binary format을 구현하지 않는다.

그냥:

\begin{dhkcode}
import numpy as np

np.save("array.npy", arr)
\end{dhkcode}

하고:

\begin{dhkcode}
arr = np.load("array.npy")
\end{dhkcode}

한다.

PyTorch도:

\begin{dhkcode}
torch.save(...)
torch.load(...)
\end{dhkcode}

이미지도:

\begin{dhkcode}
Image.open(...)
\end{dhkcode}

Pandas Parquet도:

\begin{dhkcode}
pd.read_parquet(...)
\end{dhkcode}

처럼 사용한다.

즉:

\begin{dhksmry}[centered]
저수준

open(..., 'rb')
open(..., 'wb')
\da
bytes 직접 처리



고수준

np.load()
torch.load()
Image.open()
pd.rad_parquet()
\da
라이브러리가 binary format 처리
\end{dhksmry}

그렇기에 ML/DL 하면서 \zsh{'rb'}를 거의 못 본것이 이상한 것은 아니다.

\subsubsection{File I/O 전체가 이제 하나로 연결된다}

전체 구조를 보면:

\begin{dhktree}
file system
>Path
>>open()
>>>file object
>>>>binary mode
>>>>>bytes
>>>>>>보통 전용 라이브러리
>>>>>>>NumPy
>>>>>>>Pytorch
>>>>>>>Image
>>>>text mode
>>>>>decoding
>>>>>>str
>>>>>>>TXT
>>>>>>>JSON
>>>>>>>>Dict
>>>>>>>CSV
>>>>>>>>DataFrame
\end{dhktree}

이정도면 File I/O를 거의 완료했다고 봐도 된다.

특히 중요한 건 API를 암기하는 것이 아니라:

\begin{dhkbranch}[in=8]
node main :: Path
arrow main
node main :: file object
branch main -> text :: text라면 str/encoding
arrow text
node text :: JSON, CSV 같은 구조 해석
stop text
handoff main -> binary :: binary라면 bytes
arrow binary
node binary :: 보통 전용 라이브러리가 해석
stop binary
\end{dhkbranch}

이라는 전체 레이어를 이해하는 것이다.

\begin{dhknote}
\begin{dhksmry}
    번외로, \zsh{iterator}에 대한 한 문장 설명을 추가하자면,

    Iterator는 for 문을 돌릴 수 없는 객체가 아니라, 

    for 문이 실제로 값을 하나씩 꺼내기 위해 사용하는 객체다.
\end{dhksmry}
\end{dhknote}
